\chapter{Antivirals for Influenza and COVID-19}
\index{Antivirals for Influenza and COVID-19}

\section*{Definition and Overview}
Antiviral medicines are prescription drugs that interfere with a virus's ability to replicate, helping the body clear the infection faster and reducing the risk of severe illness. This chapter covers the antivirals used against influenza (flu) and COVID-19, the illness caused by the SARS-CoV-2 coronavirus. These medicines do not cure the infection outright, and they are not a substitute for vaccination, which remains the most effective way to prevent these illnesses.

Antivirals are used mainly in people who have, or are at high risk of developing, serious illness -- for example, older adults, pregnant women, and people with chronic medical conditions or weakened immunity. They work best when started early, ideally within the first few days of symptoms. Treatment decisions are individualised according to the virus, the timing of the illness, the person's risk factors, and potential drug interactions.

\section*{Epidemiology}
Influenza causes seasonal epidemics, usually in winter in temperate regions, infecting millions of people each year and leading to thousands of hospitalisations and deaths worldwide. The burden falls mainly on the elderly, the very young, pregnant women, and people with chronic heart, lung, kidney, or metabolic disease, or weakened immunity.

COVID-19 became a global pandemic from 2020 onward and continues to circulate as an endemic respiratory virus. Its severe outcomes are concentrated in the same at-risk groups, and recurring waves have been driven by successive variants. Because both viruses evolve continually, the effectiveness of antivirals, the recommendations for their use, and the vaccine schedule can change from year to year, so current national guidance should always be consulted.

\section*{Aetiology and Pathophysiology}
Influenza is caused by the influenza A and B viruses; COVID-19 is caused by SARS-CoV-2. Both enter the body through the respiratory tract, attach to cells lining the airways, and use the cells' own machinery to replicate. Antiviral drugs interrupt this process at different steps:

\begin{itemize}
    \item \textbf{Neuraminidase inhibitors} (oseltamivir, zanamivir): block an enzyme the influenza virus needs to release new copies from infected cells and spread to neighbouring cells
    \item \textbf{Protease inhibitors} (nirmatrelvir, given with ritonavir to boost its levels): block SARS-CoV-2 from cutting its proteins into functional parts, so new virus particles cannot form
    \item \textbf{Polymerase inhibitors} (remdesivir, molnupiravir): interfere with the copying of the virus's genetic material, limiting viral load and spread
\end{itemize}

The antiviral effect is greatest when treatment begins early, before the virus has multiplied extensively -- which is why timing matters so much.

\section*{Clinical Features}
Both illnesses typically begin with fever, cough, sore throat, muscle aches, headache, and fatigue. COVID-19 frequently also causes loss or change of taste and smell, and sometimes shortness of breath, while influenza often has a more sudden onset and more prominent systemic symptoms. Most people recover at home within one to two weeks, but symptoms can be severe, particularly in at-risk groups, and a cough and tiredness may persist for weeks.

Antivirals are considered mainly in people with confirmed or suspected infection who have risk factors such as older age, pregnancy, or chronic heart, lung, kidney, liver, or immune conditions, and they are most effective when started within a few days of symptom onset. Side effects are usually mild, with nausea, vomiting, and diarrhoea the most common.

\section*{Differential Diagnosis}
\begin{itemize}
    \item \textbf{Other respiratory viruses:} respiratory syncytial virus (RSV), adenovirus, rhinovirus, and parainfluenza, which cause similar syndromes
    \item \textbf{Bacterial pneumonia:} which may follow or complicate a viral illness and requires antibiotics
    \item \textbf{Bacterial throat infections:} such as strep throat, presenting with sore throat and fever
    \item \textbf{Allergic rhinitis or asthma flares:} which cause cough, congestion, and shortness of breath without significant fever
    \item \textbf{Whooping cough (pertussis):} prolonged coughing bouts with characteristic paroxysms
    \item \textbf{Other febrile illnesses:} such as early sepsis, which can mimic a viral syndrome
\end{itemize}

\section*{When to Seek Medical Care}
See a doctor promptly if you are in a high-risk group and develop flu- or COVID-like symptoms, because starting antiviral treatment early gives the best chance of benefit. Also seek advice if symptoms are unusually severe, last longer than about a week, or worsen after initially improving, or if there is uncertainty about testing and treatment eligibility.

\textbf{Contact your local emergency services for:}
\begin{itemize}
    \item Difficulty breathing, breathlessness at rest, or persistent chest pain
    \item Blue lips or face, or confusion, drowsiness, or difficulty rousing
    \item Coughing up blood
    \item Signs of dehydration, such as little or no urine, severe dizziness, or light-headedness on standing
    \item A child who is unusually floppy, difficult to wake, or not drinking enough
    \item Seizures or a high fever with a non-blanching rash
\end{itemize}

\section*{Diagnosis and Investigations}
\begin{itemize}
    \item \textbf{Clinical assessment:} history of symptom onset, exposure to confirmed cases, vaccination status, and risk factors for severe disease
    \item \textbf{Molecular testing:} PCR or rapid antigen tests on nose or throat swabs to confirm influenza or COVID-19, and to inform treatment decisions
    \item \textbf{Blood tests:} in severe or uncertain cases, to check oxygen saturation, kidney and liver function, and inflammatory markers
    \item \textbf{Chest imaging:} chest X-ray or CT if pneumonia, effusion, or cardiac involvement is suspected
    \item \textbf{Medication review:} essential, because several antivirals -- especially nirmatrelvir--ritonavir -- interact with common medicines such as statins and anticoagulants
    \item \textbf{Monitoring in hospital:} pulse oximetry and observation for those with low oxygen levels or significant co-morbidity
\end{itemize}

\section*{Management}
\begin{itemize}
    \item \textbf{Antiviral treatment:} prescribed for at-risk people, ideally within 5 days of symptom onset for influenza (oseltamivir or zanamivir) and within 5 days for COVID-19 (nirmatrelvir--ritonavir, remdesivir, or molnupiravir), with suitability depending on age, risk, and interactions
    \item \textbf{Supportive care:} rest, fluids, and simple pain and fever relievers such as paracetamol
    \item \textbf{Oxygen and hospital care:} for people with severe disease, low oxygen levels, or significant dehydration
    \item \textbf{Management of complications:} antibiotics for secondary bacterial pneumonia and anticoagulation where clinically indicated in severe COVID-19
    \item \textbf{Vaccination:} annual influenza vaccination and recommended COVID-19 booster doses for prevention
    \item \textbf{Infection control:} staying home while infectious and reducing contact with vulnerable people
\end{itemize}

\section*{Complications}
\begin{itemize}
    \item Viral pneumonia or secondary bacterial pneumonia
    \item Worsening of asthma, heart failure, diabetes, and other chronic conditions
    \item Myocarditis, pericarditis, and other inflammatory complications
    \item Blood clots, particularly in severe COVID-19
    \item Sepsis and multi-organ failure in severe infection
    \item Side effects of antivirals, most commonly nausea, vomiting, and diarrhoea
    \item Drug interactions, especially with nirmatrelvir--ritonavir, which affects many common medicines
    \item Long COVID: persistent symptoms such as fatigue, breathlessness, and brain fog after the initial infection
\end{itemize}

\section*{Prognosis and Outcomes}
Most otherwise healthy people recover fully from influenza and COVID-19 within one to two weeks, although a cough, fatigue, and reduced exercise tolerance can persist longer. In people at high risk, antivirals started early may shorten the illness, reduce its severity, and lower the chance of hospitalisation, ventilation, and death -- particularly with nirmatrelvir--ritonavir for COVID-19.

Outcomes are worse in the elderly, the immunocompromised, and those with chronic disease, which is why early testing, vaccination, and prompt antiviral treatment matter most in these groups. Influenza and COVID-19 can both be fatal, and both can leave lasting symptoms, but the overall picture has improved substantially with vaccination and the availability of effective antiviral agents.

\section*{Prevention and Self-Care}
\begin{itemize}
    \item Keep up to date with annual influenza vaccination and recommended COVID-19 boosters
    \item Wash hands regularly, and use hand sanitiser when soap and water are unavailable
    \item Cover coughs and sneezes with a tissue or elbow, and dispose of tissues safely
    \item Improve indoor ventilation and consider masks in crowded or high-risk settings
    \item Stay home and rest while unwell to avoid spreading the virus
    \item Drink plenty of fluids and treat fever and aches with paracetamol
    \item Take antiviral medicines exactly as prescribed, and tell your doctor about all other medications and supplements
    \item Seek early review if symptoms are severe, worsening, or persisting
\end{itemize}

\section*{Sources and Further Reading}
The following organisations provide reliable, up-to-date information on influenza, COVID-19, and antiviral treatment for patients, students, and health professionals.

\begin{itemize}
    \item Centers for Disease Control and Prevention. \textit{Resources on influenza, COVID-19, and antiviral treatment.} https://www.cdc.gov/
    \item World Health Organization. \textit{Resources on influenza and COVID-19 guidance and treatment.} https://www.who.int/
    \item NHS. \textit{Information on flu, COVID-19, and antiviral medicines.} https://www.nhs.uk/conditions/
    \item NICE. \textit{Guidance relevant to COVID-19 and influenza treatment.} https://www.nice.org.uk/
    \item MedlinePlus. \textit{Health information on influenza, COVID-19, and antiviral drugs.} https://medlineplus.gov/
    \item MSD Manual. \textit{Professional overview of influenza, COVID-19, and antiviral therapy.} https://www.msdmanuals.com/
\end{itemize}
